% !TEX program = xelatex
\documentclass[12pt,a4paper]{article}

\usepackage{fontspec}
\usepackage{polyglossia}
\setmainlanguage{russian}
\setotherlanguage{english}
\setmainfont{DejaVu Serif}
\setsansfont{DejaVu Sans}
\setmonofont{DejaVu Sans Mono}

\usepackage{amsmath,amssymb,amsfonts,mathtools,bm}
\usepackage{geometry}
\usepackage{booktabs}
\usepackage{longtable}
\usepackage{array}
\usepackage{enumitem}
\usepackage{hyperref}
\usepackage{microtype}
\geometry{margin=2.5cm}
\hypersetup{colorlinks=true,linkcolor=black,urlcolor=blue,citecolor=black}
\setlength{\emergencystretch}{3em}
\sloppy

\DeclareMathOperator{\Adm}{Adm}
\DeclareMathOperator{\Struct}{Struct}
\DeclareMathOperator{\Dist}{Dist}
\DeclareMathOperator{\Dom}{Dom}
\DeclareMathOperator{\Cond}{Cond}
\DeclareMathOperator{\Loss}{Loss}
\DeclareMathOperator{\Tr}{Tr}
\DeclareMathOperator{\Diff}{Diff}
\DeclareMathOperator{\Spin}{Spin}
\DeclareMathOperator{\Var}{Var}

\newcommand{\M}{\mathcal M_U}
\newcommand{\MI}{\mathcal M_I}
\newcommand{\EI}{\mathcal E_I}
\newcommand{\Hilb}{\mathcal H}
\newcommand{\Prob}{\mathbb P}
\newcommand{\Lag}{\mathcal L}
\newcommand{\Acal}{\mathbb A_I}
\newcommand{\Fcal}{\mathbb F_I}
\newcommand{\TMI}{\mathcal T_{\mathrm{MI}}}
\newcommand{\eps}{\varepsilon}
\newcommand{\dd}{\,d}

\newcommand{\statementbox}[1]{%
\begin{center}%
\fbox{\begin{minipage}{0.9\linewidth}\centering #1\end{minipage}}%
\end{center}%
}

\title{\textbf{Теория многообразных интерфейсов}\
\large Физический слой и интерфейсно-полевая программа}
\author{Salkutsan Aleksey}
\date{Май 2026}

\begin{document}
\maketitle

\begin{abstract}
В работе представлен физический слой Теории многообразных интерфейсов как эффективная интерфейсно-полевая исследовательская программа. Центральное утверждение состоит в том, что геометрия задаёт область допустимых физических переходов, квантовая теория задаёт распределение вероятностей по этой области, а интерфейсное поле осуществляет переход от распределённых возможностей к устойчивым физическим структурам. Вводится минимальное интерфейсное поле $\chi_I$, эффективное действие, объединяющее гравитационный сектор, Стандартную модель, интерфейсный сектор, смешивание и декогеренцию, а также проверочная матрица для квантово-геометрической декогеренции, тёмного сектора, чёрных дыр и интерфейсно-полевой архитектуры объединения. Работа не заявляет завершённую квантовую гравитацию или окончательную Theory of Everything. Она формулирует структурированную физическую программу с явными условиями согласованности, проверяемыми секторами и указанными ограничениями.
\end{abstract}

\noindent\textbf{Ключевые слова:} многообразные интерфейсы; интерфейсное поле; допустимые переходы; квантовая декогеренция; тёмный сектор; чёрные дыры; объединение взаимодействий; эффективная теория поля.

\tableofcontents
\newpage

\section{Введение}
Современная фундаментальная физика содержит несколько исключительно успешных, но концептуально раздельных описаний. Общая теория относительности описывает геометрию пространства-времени через метрику $g_{\mu\nu}$ и связывает её с энергией и импульсом через уравнение Эйнштейна,
\[
G_{\mu\nu}+\Lambda g_{\mu\nu}=\frac{8\pi G}{c^4}T_{\mu\nu}.
\]
Квантовая теория описывает состояния не как заранее определённые структуры, а как суперпозиции или вероятностные распределения,
\[
|\Psi\rangle=\sum_\alpha c_\alpha |s_\alpha\rangle.
\]
Теория многообразных интерфейсов (ТМИ) предлагает общий язык для связи допустимости, вероятности и структуры. Физический слой ТМИ организован цепочкой
\[
\boxed{
g\Rightarrow \Adm_c(g)\Rightarrow \Hilb_I(g)\Rightarrow \Psi_I
\Rightarrow \Dist_{\mathrm{poss}}(\Psi_I;g)\Rightarrow \chi_I\Rightarrow \Struct_I .}
\]
Смысл этой цепочки:
\[
\boxed{\text{геометрия задаёт допустимость;}}
\qquad
\boxed{\text{квантовая теория задаёт распределение;}}
\qquad
\boxed{\text{интерфейсное поле задаёт структурирование.}}
\]

У этой работы две задачи. Первая — дать чистую математико-физическую формулировку этой цепочки. Вторая — показать, как одна и та же интерфейсно-полевая архитектура даёт четыре физических режима: квантово-геометрическую декогеренцию, феноменологию тёмного сектора, чёрные дыры как интерфейсы и интерфейсно-полевую архитектуру объединения взаимодействий.

Статус работы сознательно осторожный. ТМИ-Physics понимается как эффективная исследовательская программа. Здесь не утверждается, что уже выведена полная квантовая динамика геометрии, полная Стандартная модель или окончательная Theory of Everything. Более узкое утверждение состоит в следующем: если интерфейсное поле $\chi_I$ вводится как физический посредник между допустимыми возможностями и реализованной структурой, то можно сформулировать несколько проверяемых физических конструкций.

\section{Связь с предыдущими публикациями ТМИ}
Настоящая работа объединяет и систематизирует несколько предыдущих модулей Теории многообразных интерфейсов. В отдельных материалах рассматривались чёрные дыры, тёмный сектор, квантово-геометрическая декогеренция, математическая абстракция, осторожная квантово-гравитационная формулировка и интерфейсно-полевая архитектура объединения. В данной рукописи эти направления рассматриваются как режимы одного физического слоя. Предыдущие материалы перечислены в списке литературы и являются концептуальными предшественниками настоящего синтеза.

\section{Геометрическая допустимость}
Пусть $(\M,g)$ — пространство-время с лоренцевой метрикой. В каждой точке $x\in\M$ существует касательное пространство $T_x\M$. Определим будущий интерфейсный причинный конус:
\[
\boxed{
\mathcal C^{+}_{I,g}(x):=
\{v\in T_x\M\setminus\{0\}\mid g_x(v,v)\leq0,\ v\ \text{направлен в будущее}\}.
}
\]
Это не новый причинный конус сверх общей теории относительности. Это стандартный будущий причинный конус, интерпретированный как локальная область допустимых направлений физического перехода:
\[
\mathcal C^{+}_{I,g}(x)\equiv \mathcal C^{+}_{g}(x).
\]
Его граница задаётся условием
\[
g_x(v,v)=0,
\]
поэтому совпадает со световым конусом:
\[
\boxed{\partial\mathcal C^{+}_{I,g}(x)\equiv \mathcal C^{+}_{\mathrm{light},g}(x).}
\]
Следовательно, скорость света в ТМИ интерпретируется как граница допустимого причинного интерфейсного перехода, а не как дополнительный постулат.

Причинное будущее точки обозначается $J_g^+(x)$. Множество причинно допустимых переходов определяется так:
\[
\boxed{
\Adm_c(g):=
\{(x,y)\in\M\times\M\mid y\in J_g^+(x)\}.
}
\]
Это фиксирует первый слой теории:
\[
\boxed{g\Rightarrow \mathcal C^+_{I,g}(x)\Rightarrow J_g^+(x)\Rightarrow \Adm_c(g).}
\]
Физический смысл: геометрия задаёт не только форму пространства-времени, но и область физически допустимых переходов.

\section{Квантовое распределение по допустимой области}
Геометрия задаёт, что допустимо, но не задаёт, какой именно допустимый переход реализуется. Квантовый слой назначает амплитуды и вероятности на области допустимых переходов.

Вводится измеримое пространство
\[
(\Adm_c(g),\mathcal B_g,\mu_g),
\]
где $\mathcal B_g$ — $\sigma$-алгебра измеримых подмножеств, а $\mu_g$ — эффективная или регуляризованная мера. На фундаментальном уровне мера не считается полностью построенной. В зависимости от уровня модели она может быть дискретной мерой, минисуперпространственной мерой, цилиндрической мерой или эффективной мерой интеграла по путям.

Гильбертово пространство интерфейсных возможностей:
\[
\boxed{
\Hilb_I(g):=L^2(\Adm_c(g),\mathcal B_g,\mu_g).
}
\]
Состояние $\psi_I\in\Hilb_I(g)$ задаёт вероятностную меру
\[
\boxed{
d\Prob_{\Psi,g}(\Phi)=|\psi_I(\Phi)|^2\,d\mu_g(\Phi),
\qquad \Phi\in\Adm_c(g).
}
\]
Тем самым квантовый слой записывается как
\[
\boxed{
\Adm_c(g)\Rightarrow \Hilb_I(g)\Rightarrow \Psi_I\Rightarrow \Dist_{\mathrm{poss}}(\Psi_I;g).
}
\]
Квантовое состояние задаёт распределение по уже геометрически допустимой области.

\section{Интерфейсное структурирование}
Распределение вероятностей по допустимым переходам ещё не является реализованной структурой. ТМИ вводит интерфейсное поле, чтобы описать переход от распределённых возможностей к устойчивым физическим режимам.

В минимальной модели интерфейсное поле является вещественным скалярным полем,
\[
\boxed{\chi_I:\M\to\mathbb R.}
\]
В более общей геометрической форме оно может быть сечением интерфейсного расслоения,
\[
\boxed{\chi_I\in\Gamma(\mathcal E_\chi),\qquad \mathcal E_\chi\to\M.}
\]
Отображение структурирования записывается как
\[
\boxed{
\mathcal S_I:\Dist_{\mathrm{poss}}(\Psi_I;g)\longrightarrow \Struct_I.
}
\]
Более полно:
\[
\boxed{
\Struct_I=\mathcal S_I(g,\Psi_I,\chi_I,\Adm_c(g)).
}
\]
То есть структура возникает не из одной геометрии и не из одной волновой функции, а из совместного действия геометрической допустимости, квантового распределения и интерфейсного поля.

Первый вычислимый режим можно записать для двух геометрических альтернатив:
\[
|\Psi_g\rangle=c_1|g_1\rangle+c_2|g_2\rangle,
\qquad |c_1|^2+|c_2|^2=1.
\]
В матрице плотности недиагональный элемент
\[
\rho_{12}=c_1\overline{c_2}
\]
задаёт когерентность между двумя геометриями. Интерфейсное структурирование моделируется подавлением этой когерентности:
\[
\boxed{
|\rho_{12}(t)|=e^{-\Gamma_I t}|\rho_{12}(0)|.
}
\]
Минимальная степень структурирования:
\[
\boxed{
\Struct_I(t):=1-\frac{|\rho_{12}(t)|}{|\rho_{12}(0)|}=1-e^{-\Gamma_I t}.
}
\]

\section{Минимальное эффективное действие}
Физический слой кодируется минимальным эффективным действием
\[
\boxed{
S_{TMI}^{min}=S_{GR}+S_{SM}+S_\chi+S_{mix}+S_{dec}^{eff}.
}
\]
Развёрнуто:
\[
\boxed{
S_{TMI}^{min}
=
\int_{\M}\sqrt{-g}\left[
\frac{c^4}{16\pi G}(R-2\Lambda)
+\Lag_{SM}+\Lag_\chi+\Lag_{mix}
\right]d^4x
+S_{dec}^{eff}.
}
\]
Здесь $S_{dec}^{eff}$ — эффективный член открытой квантовой динамики, а не фундаментальное локальное действие.

Лагранжиан интерфейсного поля:
\[
\boxed{
\Lag_\chi
=-\frac12 g^{\mu\nu}\nabla_\mu\chi_I\nabla_\nu\chi_I
-V(\chi_I)-\frac12\xi_I R\chi_I^2.
}
\]
Минимальный потенциал:
\[
\boxed{
V(\chi_I)=V_0+\frac12 m_I^2\chi_I^2+\frac{\beta_I}{4}\chi_I^4.
}
\]
Физическая устойчивость требует, чтобы потенциал был ограничен снизу. Достаточным условием является $\beta_I>0$ вместе с подходящей нижней границей $V(\chi_I)$.

Чтобы не разрушать калибровочную инвариантность, $\chi_I$ берётся как синглет Стандартной модели,
\[
\chi_I\sim (1,1,0)
\quad \text{относительно}\quad SU(3)\times SU(2)\times U(1).
\]
Минимальная безопасная форма смешивания:
\[
\boxed{
\Lag_{mix}
=
-\frac14\sum_{a=1}^{3}\zeta_a\frac{\chi_I^2}{M_I^2}\Tr(F_{a\mu\nu}F_a^{\mu\nu})
+\eta_H\chi_I^2H^\dagger H
+\sum_f y_{I,f}\frac{\chi_I}{M_I}\overline{\Psi}_{L,f}H\Psi_{R,f}+h.c.
}
\]
Последний член записан в форме, совместимой с полем Хиггса, чтобы не нарушать электрослабую калибровочную инвариантность до спонтанного нарушения симметрии.

\section{Уравнения движения}
Варьирование эффективного действия даёт геометрическое уравнение
\[
\boxed{
G_{\mu\nu}+\Lambda g_{\mu\nu}
=
\frac{8\pi G}{c^4}
\left(T_{\mu\nu}^{SM}+T_{\mu\nu}^{\chi}+T_{\mu\nu}^{mix}\right)
+\mathcal D_{\mu\nu}^{eff}.
}
\]
Здесь $\mathcal D_{\mu\nu}^{eff}$ обозначает возможный эффективный вклад открытой квантовой динамики. Если обратное влияние декогеренции на геометрию не включается, можно положить $\mathcal D_{\mu\nu}^{eff}=0$.

Для минимальной скалярной связи тензор энергии-импульса интерфейсного поля имеет вид
\[
\boxed{
T_{\mu\nu}^{\chi}
=
\nabla_\mu\chi_I\nabla_\nu\chi_I
-g_{\mu\nu}\left[\frac12 g^{\alpha\beta}\nabla_\alpha\chi_I\nabla_\beta\chi_I+V(\chi_I)\right].
}
\]
Неминимальная связь даёт дополнительные члены
\[
\boxed{
T_{\mu\nu}^{\chi,\xi}
=\xi_I\left[G_{\mu\nu}\chi_I^2+g_{\mu\nu}\Box_g(\chi_I^2)-\nabla_\mu\nabla_\nu(\chi_I^2)\right].
}
\]

Варьирование по $\chi_I$ даёт
\[
\boxed{
\Box_g\chi_I-(m_I^2+\xi_I R)\chi_I-\beta_I\chi_I^3=-J_I.
}
\]
Для квантово-геометрического режима источник можно взять в виде
\[
\boxed{
J_I=\alpha_I D_g(\Psi_g),
\qquad
D_g(\Psi_g)=|c_1|^2|c_2|^2d_I^2(g_1,g_2).
}
\]
Открытая квантовая динамика записывается как
\[
\boxed{
\frac{d\rho}{dt}=-\frac{i}{\hbar}[\widehat H_{eff},\rho]+\mathcal L_{dec}^{I}[\rho].
}
\]
Для двух геометрий получаем
\[
\boxed{|\rho_{12}(t)|=e^{-\Gamma_I t}|\rho_{12}(0)|.}
\]

\section{Режим квантово-геометрической декогеренции}
Квантово-геометрический режим возникает, когда квантовое состояние содержит суперпозицию геометрических альтернатив,
\[
|\Psi_g\rangle=c_1|g_1\rangle+c_2|g_2\rangle.
\]
Геометрический контраст:
\[
\boxed{
D_g(\Psi_g)=|c_1|^2|c_2|^2d_I^2(g_1,g_2).
}
\]
Интерфейсный источник возбуждает $\chi_I$, а интерфейсное поле создаёт скорость декогеренции:
\[
\boxed{
\Psi_g\Rightarrow D_g\Rightarrow J_I\Rightarrow \chi_I\Rightarrow \Gamma_I\Rightarrow \Struct_I.
}
\]
Минимальная феноменологическая форма:
\[
\boxed{
\Gamma_I=\kappa_I|c_1|^2|c_2|^2d_I^2(g_1,g_2).
}
\]
В минисуперпространственной модели с масштабным фактором $a$:
\[
\boxed{
d_I(g(a_1),g(a_2))=\left|\ln\frac{a_1}{a_2}\right|,
}
\]
следовательно
\[
\boxed{
\Gamma_I=\kappa_I|c_1|^2|c_2|^2\ln^2\left(\frac{a_1}{a_2}\right).
}
\]
Наблюдаемая избыточная декогеренция:
\[
\boxed{
\Delta\Gamma_I=\Gamma_{obs}-\Gamma_{std}=\kappa_I x_I,
\qquad
x_I=|c_1|^2|c_2|^2\ln^2\left(\frac{a_1}{a_2}\right).
}
\]
Нулевая гипотеза: $H_0^{QG}:\kappa_I=0$. Гипотеза ТМИ: $H_{TMI}^{QG}:\kappa_I>0$.

\section{Режим тёмного сектора}
Режим тёмного сектора основан на гипотезе
\[
\boxed{
T_{\mu\nu}^{dark}\sim T_{\mu\nu}^{\chi}.
}
\]
Поле раскладывается на глобальную и локальную компоненты:
\[
\boxed{\chi_I(x,t)=\chi_0(t)+\delta\chi(x,t).}
\]
Глобальная компонента связывается с режимом тёмной энергии,
\[
\boxed{\chi_0(t)\Rightarrow DE,}
\]
локальная компонента — с режимом тёмной материи,
\[
\boxed{\delta\chi(x,t)\Rightarrow DM.}
\]
Для фоновой компоненты
\[
\boxed{
\rho_{DE}^{I}=\frac12\dot\chi_0^2+V(\chi_0),
\qquad
p_{DE}^{I}=\frac12\dot\chi_0^2-V(\chi_0),
}
\]
поэтому
\[
\boxed{
w_I=\frac{p_{DE}^{I}}{\rho_{DE}^{I}}
=\frac{\frac12\dot\chi_0^2-V(\chi_0)}{\frac12\dot\chi_0^2+V(\chi_0)}.
}
\]
Если $\dot\chi_0^2\ll V(\chi_0)$, то $w_I\approx -1$.

Для локальной компоненты приближённое уравнение моды:
\[
\boxed{
\delta\ddot\chi+3H\delta\dot\chi-\frac{\nabla^2}{a^2}\delta\chi+m_{eff}^2\delta\chi=0,
}
\]
где
\[
\boxed{m_{eff}^2=m_I^2+3\beta_I\chi_0^2+\xi_I R.}
\]
Осциллирующие локальные моды могут вести себя как бездавленческая компонента: $\langle p_{DM}^{I}\rangle\approx0$ и $\rho_{DM}^{I}\propto a^{-3}$.

Проверка тёмного сектора не означает, что ТМИ уже доказала природу реальной тёмной материи и тёмной энергии. Условная проверка такова: если тёмный сектор создаётся полем $\chi_I$, то возможные отклонения должны появиться в наборе
\[
\boxed{
\mathcal O_{Dark}=\{w_I(z),H(z),\Delta v_c^2(r),\Delta\alpha_{lens},G_{eff}(k,z),\Delta\Gamma_I^{dark}\}.
}
\]

\section{Режим чёрной дыры как интерфейса}
Чёрная дыра интерпретируется как предельный причинно-информационный интерфейс. В стандартных геометрических терминах
\[
\boxed{
\mathcal B_g=\M\setminus J_g^-(\mathcal I^+),
\qquad
\mathcal H_g=\partial\mathcal B_g.
}
\]
В ТМИ горизонт событий является границей внешней допустимости:
\[
\boxed{\mathcal H_g=\partial\Adm_{out}.}
\]
Для геометрии Шварцшильда
\[
\boxed{
ds^2=-\left(1-\frac{r_s}{r}\right)c^2dt^2+\left(1-\frac{r_s}{r}\right)^{-1}dr^2+r^2d\Omega^2,
\qquad
r_s=\frac{2GM}{c^2}.
}
\]
Радиальные световые лучи удовлетворяют
\[
\boxed{\frac{dr}{dt}=\pm c\left(1-\frac{r_s}{r}\right).}
\]
Горизонт отмечает границу, где классический внешний канал закрывается. В интерфейсной записи
\[
\boxed{r<r_s\Rightarrow \Adm_{out}^{class}(I_{BH})=\varnothing,
\qquad \Adm_{in}(I_{BH})\neq\varnothing.}
\]

Энтропия горизонта интерпретируется как информационная ёмкость интерфейса:
\[
\boxed{
S_{BH}=\frac{k_Bc^3A}{4G\hbar}=\mathrm{Cap}_{info}(I_{BH}).
}
\]
Число битов:
\[
\boxed{N_{bits}=\frac{A}{4l_P^2\ln2}.}
\]
Возможная наблюдаемая величина TMI-BH — нетепловые корреляции в излучении Хокинга:
\[
\boxed{
\rho_{rad}=\rho_{thermal}+\rho_{corr},
\qquad
\Cond_{info}^{quant}>0\Rightarrow \rho_{corr}\neq0.
}
\]
Также возможна гипотеза усиленной декогеренции около чёрной дыры:
\[
\boxed{
\Gamma_{BH}^{I}=\gamma_I\frac{S_{BH}}{k_B}.
}
\]
Это феноменологическая гипотеза, а не вывод полной квантовой гравитации.

\section{Интерфейсно-полевая архитектура объединения}
Раздел объединения формулируется как интерфейсно-полевая архитектура, а не как завершённая Theory of Everything. Четыре фундаментальных взаимодействия рассматриваются как интерфейсные каналы:
\[
\boxed{
\mathcal F_{phys}=\{I_G,I_{EM},I_W,I_S\}.
}
\]
Они интерпретируются как проекции универсальной физической интерфейсной структуры.

Вводится универсальное интерфейсное расслоение
\[
\boxed{\pi_I:\mathcal E_I\to\M,}
\]
и универсальная интерфейсная связность
\[
\boxed{\Acal\in\Omega^1(\M,\mathfrak g_I).}
\]
Безопасное условие включения:
\[
\boxed{G_I\supset \Spin(1,3)\times SU(3)\times SU(2)\times U(1).}
\]
В минимальной архитектуре можно записать
\[
\boxed{
G_I=G_{grav}\ltimes G_{SM}\ltimes G_\chi,
\qquad
G_{SM}=SU(3)\times SU(2)\times U(1).
}
\]
Связность раскладывается так:
\[
\boxed{
\Acal=\omega^{ab}J_{ab}+e^aP_a+G^AT_A+W^iT_i+BY+A_\chi\Xi.
}
\]
Кривизна:
\[
\boxed{\Fcal=d\Acal+\Acal\wedge\Acal.}
\]
Секторные проекции являются отображениями
\[
\boxed{\pi_a:\mathfrak g_I\to\mathfrak g_a,
\qquad a\in\{G,S,W,EM\},}
\]
при этом
\[
\boxed{F_a=\pi_a(\Fcal).}
\]
Первый феноменологический выход — интерфейсная поправка к калибровочным связям:
\[
\boxed{
\frac{1}{g_{a,eff}^2}=\frac{1}{g_a^2}+\zeta_a\frac{\chi_I^2}{M_I^2},
\qquad
\Delta_a(E)=\zeta_a\frac{\chi_I^2(E)}{M_I^2}.
}
\]

\section{Единая наблюдаемая матрица}
Физический слой является проверяемым только тогда, когда его режимы связаны с наблюдаемыми величинами. Предлагаемая матрица:
\[
\boxed{
\mathcal T_{TMI}^{phys}
=
\left(
\Delta\Gamma_I,
 w_I(z),
 G_{eff}(k,z),
 \rho_{corr},
 \Delta_a(E)
\right).
}
\]
Смысл:
\begin{longtable}{p{0.25\linewidth}p{0.32\linewidth}p{0.32\linewidth}}
\toprule
Режим & Наблюдаемая величина & Нулевая гипотеза \\
\midrule
QG & $\Delta\Gamma_I$ & нет избыточной геометрической декогеренции \\
Dark & $w_I(z),G_{eff}(k,z)$ & $\Lambda$CDM-подобная независимость тёмной материи и тёмной энергии \\
BH & $\rho_{corr}$ & чисто тепловое излучение в соответствующем приближении \\
ToE & $\Delta_a(E)$ & нет интерфейсной поправки к калибровочным связям \\
\bottomrule
\end{longtable}

Статистическая форма QG-модуля:
\[
\boxed{
y_i=\Delta\Gamma^{(i)}=\kappa_Ix_I^{(i)}+\eps_i,
\qquad
\mathbb E[\eps_i]=0.
}
\]
При равных дисперсиях
\[
\boxed{
\widehat\kappa_I=\frac{\sum_i x_I^{(i)}y_i}{\sum_i(x_I^{(i)})^2}.
}
\]
Минимальная QG-гипотеза получает поддержку только если $\widehat\kappa_I>0$ и доверительный интервал не содержит ноль. Она ослабляется или отклоняется, если $\Delta\Gamma_I$ не масштабируется с $x_I$.

\section{Условия согласованности}
Физический слой должен восстанавливать известные теории в соответствующих пределах:
\[
\boxed{\chi_I\to0,\ S_{mix}\to0\Rightarrow TMI\to GR+SM.}
\]
Если убрать материю и калибровочные поля, предел должен сводиться к общей теории относительности:
\[
\boxed{S_{matter}\to0\Rightarrow TMI\to GR.}
\]
При фиксированной геометрии и исчезающей интерфейсной декогеренции
\[
\boxed{\chi_I\to0,\ \Gamma_I\to0\Rightarrow \text{стандартная квантовая эволюция}.}
\]
Для сектора объединения
\[
\boxed{E\ll M_I,
\quad \chi_I\to0,
\quad \zeta_a,\eta_H,y_{I,f}\to0
\Rightarrow S_{TMI}\to S_{SM}.}
\]
Эти условия являются обязательными. Без них интерфейсная модель была бы несовместима с установленной физикой.

\section{Уровни утверждений и ограничения}
Необходимо различать три уровня утверждений.

\subsection*{Уровень A: внутренние следствия}
Это утверждения, следующие из определений. Например,
\[
g\Rightarrow\Adm_c(g)
\]
следует после определения $\Adm_c(g)$ через $J_g^+(x)$. Аналогично,
\[
\psi_I\in L^2(\Adm_c(g),\mathcal B_g,\mu_g)
\Rightarrow
 d\Prob_{\Psi,g}=|\psi_I|^2d\mu_g.
\]

\subsection*{Уровень B: феноменологические гипотезы}
Это проверяемые физические предположения, например
\[
\Gamma_I=\kappa_I|c_1|^2|c_2|^2d_I^2(g_1,g_2),
\qquad
T_{\mu\nu}^{dark}\sim T_{\mu\nu}^{\chi}.
\]
Их нельзя считать установленными фактами природы.

\subsection*{Уровень C: открытые задачи}
В данной работе не решается полная проблема меры на пространстве геометрий, не выводится $\kappa_I$ из фундаментальных констант, не выводится полная Стандартная модель, не объясняются поколения фермионов и не строится полная квантовая теория метрики. Это исследовательские задачи, а не скрытые допущения.

\section{Заключение}
Физический слой Теории многообразных интерфейсов сводится к цепочке
\[
\boxed{
g\Rightarrow\Adm_c(g)\Rightarrow\Hilb_I(g)\Rightarrow\Psi_I
\Rightarrow\Dist_{\mathrm{poss}}(\Psi_I;g)
\Rightarrow\chi_I\Rightarrow\Struct_I.
}
\]
Теория утверждает: геометрия задаёт допустимость, квантовая теория задаёт распределение, а интерфейсное поле задаёт структурирование. Минимальная физическая модель записывается как
\[
\boxed{
S_{TMI}^{min}=S_{GR}+S_{SM}+S_\chi+S_{mix}+S_{dec}^{eff}.
}
\]
В этой рамке QG, Dark, BH и ToE являются не независимыми метафорами, а режимами одной интерфейсно-полевой архитектуры:
\[
\boxed{(g_{\mu\nu},\Psi,\chi_I,\Acal)\Rightarrow(QG,Dark,BH,ToE).}
\]
Результат не является финальной фундаментальной теорией. Это структурированная эффективная исследовательская программа с явными определениями, условиями согласованности, проверяемыми наблюдаемыми величинами и открытыми задачами.

\appendix
\section{Реестр обозначений}
\begin{longtable}{p{0.25\linewidth}p{0.65\linewidth}}
\toprule
Обозначение & Значение \\
\midrule
$\M$ & пространство-время Вселенной \\
$g_{\mu\nu}$ & лоренцева метрика \\
$\mathcal C^+_{I,g}(x)$ & будущий интерфейсный причинный конус \\
$J_g^+(x)$ & причинное будущее точки $x$ \\
$\Adm_c(g)$ & множество причинно допустимых переходов \\
$\Hilb_I(g)$ & гильбертово пространство интерфейсных возможностей \\
$\mu_g$ & эффективная или регуляризованная мера на допустимых переходах \\
$\Psi_I,\psi_I$ & интерфейсное квантовое состояние и волновая функция \\
$\Dist_{\mathrm{poss}}$ & распределение возможностей \\
$\chi_I$ & интерфейсное поле \\
$\Gamma_I$ & скорость интерфейсной декогеренции \\
$\Struct_I$ & структурированный физический режим \\
$\Acal$ & универсальная интерфейсная связность \\
$\Fcal$ & универсальная интерфейсная кривизна \\
$\Delta\Gamma_I$ & избыточная интерфейсная декогеренция \\
$w_I(z)$ & параметр состояния интерфейсной тёмной энергии \\
$\rho_{corr}$ & корреляционный вклад в излучение чёрной дыры \\
$\Delta_a(E)$ & интерфейсная поправка к калибровочным связям \\
\bottomrule
\end{longtable}

\section{Логическое ядро}
Определения задают $\M$, $g$, $\Adm_c(g)$, $\Hilb_I(g)$, $\Prob_{\Psi,g}$, $\chi_I$ и $\Struct_I$. Основные постулаты:
\[
\boxed{P_1:\ g\Rightarrow\Adm_c(g),}
\qquad
\boxed{P_2:\ \Psi_I\Rightarrow\Dist_{\mathrm{poss}}(\Psi_I;g),}
\qquad
\boxed{P_3:\ \chi_I\Rightarrow\Struct_I.}
\]
Феноменологические гипотезы — это конкретные законы связи с измеримыми величинами: $\Delta\Gamma_I=\kappa_Ix_I$, $T_{\mu\nu}^{dark}\sim T_{\mu\nu}^{\chi}$, $\rho_{rad}=\rho_{thermal}+\rho_{corr}$ и $\Delta_a(E)=\zeta_a\chi_I^2(E)/M_I^2$.

\section{Открытые исследовательские задачи}
Ближайшие технические задачи: построить хотя бы одну полностью явную меру $\mu_g$ в вычислимой модели; вывести или ограничить $\kappa_I$ из действия интерфейсного поля; сравнить законы декогеренции $p=1$ и $p=2$,
\[
\Delta\Gamma_I=\kappa_Ix_I^p,
\qquad p\in\{1,2\};
\]
превратить тёмный сектор в численную космологическую модель; уточнить критерий корреляций для чёрных дыр; определить группу объединения $G_I$ строже, чем через безопасное условие включения.

\begin{thebibliography}{99}
\bibitem{TMI-BH} Salkutsan Aleksey, \textit{Black Holes in the Theory of Manifold Interfaces}, Zenodo. \href{https://doi.org/10.5281/zenodo.20025226}{https://doi.org/10.5281/zenodo.20025226}.
\bibitem{TMI-ToE} Salkutsan Aleksey, \textit{Interface-Field Architecture for the Unification of Fundamental Interactions}, Zenodo. \href{https://doi.org/10.5281/zenodo.20025731}{https://doi.org/10.5281/zenodo.20025731}.
\bibitem{TMI-Dark} Salkutsan Aleksey, \textit{The Dark Sector as an Interface Field. A Module of the Theory of Manifold Interfaces}, Zenodo. \href{https://doi.org/10.5281/zenodo.20025682}{https://doi.org/10.5281/zenodo.20025682}.
\bibitem{TMI-QG20} Salkutsan Aleksey, \textit{Theory of Manifold Interfaces TMI-QG-2.0: A Minimal Predictive Model of Interface-Induced Geometric Decoherence}, Zenodo. \href{https://doi.org/10.5281/zenodo.20013976}{https://doi.org/10.5281/zenodo.20013976}.
\bibitem{TMI-Math} Salkutsan Aleksey, \textit{Mathematical Abstraction of the Theory of Manifold Interfaces}, Zenodo. \href{https://doi.org/10.5281/zenodo.19978895}{https://doi.org/10.5281/zenodo.19978895}.
\bibitem{TMI-CautiousQG} Salkutsan Aleksey, \textit{Theory of Manifold Interfaces. A Cautious Interface Approach to Quantum Gravity}, Zenodo. \href{https://doi.org/10.5281/zenodo.19988505}{https://doi.org/10.5281/zenodo.19988505}.
\bibitem{Einstein1915} A. Einstein, \textit{Die Feldgleichungen der Gravitation}, Sitzungsberichte der Preussischen Akademie der Wissenschaften, 1915.
\bibitem{Wald1984} R. M. Wald, \textit{General Relativity}, University of Chicago Press, 1984.
\bibitem{Bekenstein1973} J. D. Bekenstein, \textit{Black holes and entropy}, Physical Review D, 7(8), 1973.
\bibitem{Hawking1975} S. W. Hawking, \textit{Particle creation by black holes}, Communications in Mathematical Physics, 43, 1975.
\bibitem{Zurek2003} W. H. Zurek, \textit{Decoherence, einselection, and the quantum origins of the classical}, Reviews of Modern Physics, 75, 2003.
\bibitem{Schlosshauer2007} M. Schlosshauer, \textit{Decoherence and the Quantum-to-Classical Transition}, Springer, 2007.
\bibitem{Weinberg1995} S. Weinberg, \textit{The Quantum Theory of Fields}, Vol. 1, Cambridge University Press, 1995.
\end{thebibliography}

\end{document}
